\scanpage{11}
\deleted{notant que}

en particulier, tenant compte de $3^\circ$,
\[
I^{\unsure{\alpha|*|+}}(G)\longrightarrow \mathrm{Pol}(G)\times\mathrm{Eff}(G)
\]
\margin{\unsure{''} les $i:G_{m\,C}\to G_C$ définissant une polarisation, modulo \unsure{action} de $G(\mathbb R)$ \ldots}

\uline{Questions} Cette dernière application, ou $I(G)\unsure{-|\to}\mathrm{Eff}(G)$, injectives ?

Les structures $\in\mathrm{Eff}(G)$ provenant de la théorie des motifs peuvent elles tjrs se décrire par un $i_1$ ? i.e. sont elles dans $\operatorname{Im}I(G)$ ?
\margin{N.B. Le fait qu'une structure d'effectivité puisse se décrire par un élément de $I(G)$ resp $I^*(G)$ est une propriété intrinsèque à la catégorie de motifs $M$, indép. du foncteur qui la réalise. ! !}

(Regarder ce qu'on obtient en tordant par un torseur). Examiner aussi, pour ces questions et celles posées dans $3^\circ$, le cas où $G$ avec ses 4 structures \unsure{provient|provenant} d'un motif sur schéma \inserted{connexe} $S$ muni d'un $\xi\in S(\mathbb{Q})$, en utilisant le foncteur de De-Rham \ldots\ ? ]

Regarder $G$ commutatif ! $G=\unsure{\mathrm{Sl}|\mathrm{SL}|\mathrm{Gal}}(\unsure{2|n|?|\overline{\mathbb Q}},\mathbb{Q})$ !

Des propriétés arithmétiques plus profondes pourront être sans doute formulées par l'introduction des éléments de \unsure{frobenius|Frobenius|filtrations} et l'interprétation des structures $3^\circ$ et $4^\circ$ en termes de ces derniers.
% archive sheet number: 10
